% Generated draft: requires mathematical and source review.
\documentclass[11pt]{article}
\usepackage[margin=25mm]{geometry}
\usepackage{fontspec}
\setmainfont{Latin Modern Roman}
\usepackage{amsmath,amssymb,mathtools,unicode-math}
\setmathfont{Latin Modern Math}
\usepackage{enumitem,needspace,xurl,hyperref}
\setlength{\parindent}{0pt}
\setlength{\parskip}{5pt}
\setlength{\emergencystretch}{3em}
\newcommand{\N}{\mathbb{N}}
\newcommand{\Z}{\mathbb{Z}}
\newcommand{\Q}{\mathbb{Q}}
\newcommand{\R}{\mathbb{R}}
\newcommand{\C}{\mathbb{C}}
\newcommand{\F}{\mathbb{F}}
\newcommand{\A}{\mathbb{A}}
\newcommand{\PP}{\mathbb{P}}
\newcommand{\E}{\mathbb{E}}
\newcommand{\eps}{\varepsilon}
\newcommand{\dd}{\,\mathrm{d}}
\begin{document}
\section*{UnsolvedMath ID: 11000172}
\textbf{Source problem code:} catalogue entry
\subsection{Definitions and mathematical statement}
Here the source's \(M_g\) denotes a mapping class \emph{group}, not the moduli space. Write it as \(G_g\), with homology representation
\[\pi:G_g\to\operatorname{Sp}(2g,\Z),\qquad\mathcal I_g=\ker\pi.\]
The programme is to study finite quotients \(\phi:G_g\twoheadrightarrow Q\) that do not factor through \(\pi\). When \(\pi\) is the standard surjective homology representation, this means there is no \(\overline\phi\) with \(\phi=\overline\phi\circ\pi\), equivalently
\[\mathcal I_g\not\subseteq\ker\phi.\]
The task is a systematic computer-assisted investigation and organization of such quotients, not merely proof that one exists.

The source proposes starting perhaps at \(g=3\) and enumerating homomorphisms \(G_g\to\Sigma_n\) for increasing \(n\). Given a finite presentation, assign permutations to every generator and test all defining relations. Each successful assignment defines a homomorphism and its image is a finite quotient; surjectivity onto all of \(\Sigma_n\) is not required. A catalogue must specify its duplication/equivalence convention, for example conjugacy of permutation representations or isomorphism of image quotients, neither fixed in the excerpt.

The source points to [52] for a simple presentation of \(M_{g,0,0}\). It reports that factorization through the symplectic quotient is tested using two additional relation types: the genus-one separating-curve twist for \(g\geq2\), and the genus-one bounding-pair relation for \(g\geq3\), citing [23]. In conventional terms the second is a bounding-pair map \(t_a t_b^{-1}\) for homologous disjoint curves bounding the appropriate subsurface; the extracted phrase ``Dehn twist on a bounding pair'' needs that convention. The exact normal-generation statement and test words must be verified before treating these as a sufficient computational test.

Further necessary restrictions are reported: generators chosen as twists in nonseparating curves must have conjugate images, and [7] gives generation by six involutions, with its genus hypotheses not reproduced. These restrictions can reduce enumeration but do not alone ensure that a tuple defines a homomorphism. The claim that modern computers put the programme within reach is historical motivation, not a supplied bound or completed catalogue. The ending announces a related nonfinite-quotient problem that is absent.

The research summary reports existence of nonsymplectic finite quotients but only an incomplete ``\(2^{?}\)'' genus-three example and no usable reference. Such existence, even if verified, would not settle the systematic study requested here.
\subsection{Short English statement}
Systematically enumerate and understand finite mapping-class quotients that detect Torelli elements, using permutation representations and relation tests. The source asks for an organized investigation, so an isolated existence result does not finish the programme.
\subsection{Sources}
\begin{itemize}
\item Local UnsolvedMath registry, permanent ID \(11000172\), exact source code \texttt{catalogue entry}.
\item Farb, \emph{Problems on Mapping Class Groups and Related Topics}, Problem 3.1, PDF page 160: \url{https://www.math.uchicago.edu/~farb/papers/mcgbook.pdf}. Registry access 2026-07-29; pdf-font-structure extraction.
\item Internal references [52] (presentation), [23] (symplectic-factorization relations) and [7] (six involutions) lack full entries. The summary's unspecified genus-three quotient has no exact order or primary citation.
\end{itemize}
\end{document}
